\documentclass[10pt,a4paper]{amsart}
\usepackage[margin=19mm]{geometry}
\usepackage{amsmath,amssymb,mathtools}
\usepackage[scr=rsfs,cal=euler]{mathalfa}
\usepackage{xfrac}
\usepackage[unicode,hidelinks,bookmarks=false]{hyperref}
\newcommand{\ie}{i.e.\ }
\newcommand{\cf}{cf.\ }
\newcommand{\resp}{respectively\ }
\newcommand{\Subsection}{\subsection}
\providecommand{\nobreakdash}{\nobreak\hbox{-}}
\setlength{\emergencystretch}{2em}
\multlinegap0pt
\usepackage{stackengine}
\usepackage{enumerate}
\usepackage{mathtools}
\usepackage{amssymb}
\newtheorem{proposition}{Proposition}
\newtheorem{lemma}{Lemma}
\DeclareMathOperator{\gap}{{\mathsf{gap}}}
\DeclareMathOperator{\zcl}{{\mathsf{zcl}}}
\title{On the sequential topological complexity and the LS-category of the cofiber of higher diagonals for symmetric products of non-orientable surfaces}
\author{Jes\'us Gonz\'alez, Ekansh Jauhari}
\date{Source: arXiv:2603.14788v1 (CC BY 4.0)}
\begin{document}
\maketitle

\noindent\emph{Source and licence.} On the sequential topological complexity and the LS-category of the cofiber of higher diagonals for symmetric products of non-orientable surfaces. Version arXiv:2603.14788v1 (CC BY 4.0), distributed by arXiv under the Creative Commons Attribution 4.0 International licence, http://creativecommons.org/licenses/by/4.0/, as recorded in the arXiv licence metadata for this exact version and shown on the abstract page https://arxiv.org/abs/2603.14788v1. Full licence notice: \texttt{licenses/arxiv-2603.14788-CC-BY-4.0.txt}.\par\smallskip

\noindent\textit{Editorial scope.} Counted unit: the article's lemma final (source0.tex lines 1292--1299). The article's complete original proof of it is delivered with it (source0.tex lines 1303--1365, one \texttt{proof} environment of 63 lines, which the article places away from the statement). No mathematics is written, repaired or re-derived: the statement and the proof are the article's own lines, in the article's own notation, and no retained line is substituted or re-flowed.  Retained uncounted as the source-local context the proof needs: lines 869--874; lines 882--892; lines 1273--1275.  Nothing was omitted from any retained span, so the package carries no omission record.  No retained line contains a citation, so the wrapper carries no bibliography and no bibliography entry.  No retained line contains a figure environment, an image-inclusion command or drawing code, so no image is delivered and the package declares no asset dependency.  Editorial formatting only, in addition: an AMS \texttt{amsart} preamble that restates the article's own declarations and macros used by the retained lines, namely the environments \texttt{proposition}, \texttt{lemma} on separate counters, together with the standard packages those lines need.  The retained environments print in this document as Proposition 1, Lemma 1 in order of appearance, whereas the article prints the same items with the numbering of its own full document.  The complete native source of the article as distributed in the arXiv e-print for this exact version is bundled unchanged as \texttt{sources/196432/source0.tex}.
\begin{proposition}\label{prop: prop y from notes}
Let $p_{r}\in\{0,\ldots, 2n\}$ for $1\le r \le k$. If $\sum_{r=1}^kp_r\le A$, then
\[
c\left(x_{1,1}^{2n-p_1}\otimes\cdots\otimes x_{k,1}^{2n-p_k}\right)\le\left\lfloor\frac{A}{2}\right\rfloor.
\]
\end{proposition}

\begin{proposition}\label{prop: prop z from notes}
Let $b_1\otimes\cdots\otimes b_k\in H^*(SP^n(N_g))^{\otimes k}\hookrightarrow H^*(SP^n(N_{g+1}))^{\otimes k}$ be a basis element. Then for any basis element $\alpha$ in the expansion of a product
\[
p=b_1\otimes\cdots\otimes b_k\cdot\prod_{r=2}^k\left(x_{1,g+1}+x_{r,g+1}\right)^{\varepsilon_r},
\]
we have the inequality
\[
c(\alpha)\le c\left(b_1\otimes\cdots\otimes b_k\right)-\left\lfloor\frac{1+\sum_{r=2}^k\varepsilon_r}{2}\right\rfloor.
\]
In particular, the product $p$ vanishes if $c(b_1\otimes\cdots\otimes b_k)<\left\lfloor\tfrac{1+\sum_{r=2}^k\varepsilon_r}{2}\right\rfloor$.
\end{proposition}

\begin{align} 
m&=\prod_{r=2}^k\left(x_{1,1}+x_{r,1}\right)^{\alpha_{r,1}} \in H^*(P^{2n})^{\otimes k}, \ \text{ with $\alpha_{r,1}\geq0$},\ \text{ and} \nonumber \\
e&=\prod_{\substack{2\le r\le k\\ 2\le i \le g}}\left(x_{1,i}+x_{r,i}\right)^{\delta_{r,i}} \in H^*(SP^n(N_g))^{\otimes k}, \ \text{ with $\delta_{r,i}\in\{0,1\}$}, \label{eq: e0}   \end{align}

\begin{lemma}\label{prop: final}
    Let $n,g\geq1$, $d\geq0$, $k=2\lambda\ge 4$, and $\gap_k(P^{2n}):=2h-1\geq1$ such that $(g-1)(\lambda-1)<h\le(g-1)\lambda$. Let $t:=h-(g-1)(\lambda-1)$ and assume that
\begin{enumerate}
        \item $m\in H^*(P^{2n})^{\otimes k}$ is homogeneous with $\deg(m)=\zcl_k(P^{2n})-d$, and\hspace{.5mm}\footnote{We note that the equality $\deg(m)=\zcl_k(P^{2n})-d$ makes sense since the hypotheses of Lemma~\ref{prop: final} imply $d<\gap_k(P^{2n})$, whose verification is easy and hence omitted.}
        \item  $e\in H^*(SP^n(N_g))^{\otimes k}$ is as in~\eqref{eq: e0} with $\deg(e)>\gap_k(P^{2n})+d-t$.
    \end{enumerate}
    Then the product $m\cdot e$ vanishes in $H^*(SP^n(N_g))^{\otimes k}$.
\end{lemma}

\begin{proof}[Proof of Lemma~\ref{prop: final}]
We can safely assume $m=x_{1,1}^{2n-p_1}\otimes\cdots\otimes x_{k,1}^{2n-p_k}$, with $p_r\ge 0$ for $1\le r\le k$, as well as $\deg(e)=\gap_k(P^{2n})+d-t+1$, which, by hypothesis, agrees with $d+h+(g-1)(\lambda-1)$. Then $p_1+\cdots+p_k=2h-1+d$, so that Proposition~\ref{prop: prop y from notes} gives
\[
c(m)\le\left\lfloor\frac{2h-1+d}{2}\right\rfloor=h+\left\lfloor\frac{d-1}{2}\right\rfloor.
\]
The vanishing of $m\cdot e$ then follows from an iterated application of Proposition~\ref{prop: prop z from notes} and the fact that
\[
\sum_{i=2}^g\left\lfloor \frac{\delta_i+1}{2} \right\rfloor>h+\left\lfloor\frac{d-1}{2}\right\rfloor,
\]
where $\delta_i:=\sum_{r=2}^k\delta_{r,i}$.  In turn, the previous inequality is a special case of the purely algebraic relation
\begin{equation}\label{alternativa}
\min\left\{ \,\sum_{i=2}^g\left\lfloor \frac{\varepsilon_i+1}{2} \right\rfloor\ \middle|\ 
\begin{matrix}
0\le\varepsilon_i\le k-1, &  \\ \displaystyle\sum_{i=2}^g\varepsilon_i=d+h+(g-1)(\lambda-1) &
\end{matrix}\hspace{-3mm}\right\}>h+\left\lfloor\frac{d-1}2\right\rfloor,
\end{equation}
whose validity is argued next.

First, note that $\delta_i\leq k-1=2\lambda-1$ for $2\leq i\leq g$, so 
\[
(g-1)(2\lambda-1)\geq\sum_{i=2}^g\delta_i=\deg(e)=d+h+(g-1)(\lambda-1),
\]
which yields the third inequality in
$(g-1)(\lambda-1)<h\leq d+h\le (g-1)\lambda$
%To prove $m\cdot e=0$, it suffices to show, in view of the iterated application of Proposition~\ref{prop: prop z from notes}, the strict inequality in
%\begin{equation}\label{eq: main thing to prove}
%\sum_{i=2}^g\left\lfloor \frac{\delta_i+1}{2} \right\rfloor\ge\varepsilon>h+\left\lfloor\frac{d-1}{2}\right\rfloor\ge c(m).  
%\end{equation}
(the other two hold by hypothesis). We can then write $d+h=(g-1)\lambda-j$ for some $j\in\{0,\ldots,g-2\}$, so that every summation $\sum\varepsilon_i$ in~\eqref{alternativa} equals $(g-1)(k-1)-j$. It follows that any tuple $\varepsilon:=(\varepsilon_2,\ldots,\varepsilon_g)$ having $j$ coordinates equal to $k-2$ and the rest equal to $k-1$, which is unique up to permutation, gives an instance of the set in~\eqref{alternativa} that also satisfies
\begin{align*}
\sum_{i=2}^g\left\lfloor \frac{\varepsilon_i+1}{2} \right\rfloor &= \sum_{i=2}^{j+1}\left\lfloor \frac{k-1}{2}\right\rfloor+\sum_{i=j+2}^g\left\lfloor \frac{k}{2}  \right\rfloor\\&=j(\lambda-1)
+(g-j-1)\lambda=\lambda(g-1)-j>h+\left\lfloor\frac{d-1}2\right\rfloor,
\end{align*}
where the last inequality holds because $h=(g-1)\lambda-j-d$ and $d>\left\lfloor\frac{d-1}2\right\rfloor$. The inequality in~\eqref{alternativa} will then follow once we argue that
\begin{equation}\label{serealiza}
\sum_{i=2}^g\left\lfloor\frac{\varepsilon'_i+1}2\right\rfloor\geq \sum_{i=2}^g\left\lfloor\frac{\varepsilon_i+1}2\right\rfloor
\end{equation}
holds for any other instance $\varepsilon':=(\varepsilon'_2,\ldots,\varepsilon'_g)$ of the set in~\eqref{alternativa}. For this, note that, just as the tuple $\varepsilon$ has \emph{exactly} $g-j-1$ coordinates (say the first ones) with the maximal possible value $k-1$, the tuple $\varepsilon'$ must have \emph{at least} $g-j-1$ coordinates (say again the first ones) with the maximal possible value $k-1$. We can ignore these first $g-j-1$ coordinates, as they represent no difference for both sides in~\eqref{serealiza}. From the remaining coordinates, we can likewise ignore those in $\varepsilon'$ together with the corresponding ones in $\varepsilon$ with value $k-2$. Let $R\subset\{2,\ldots,g\}$ denote the set of the remaining coordinates. All of these then have value $k-2$ in the case of $\varepsilon$, while value different from $k-2$ in the case of $\varepsilon'$. Let $R'\subseteq R$ (resp.~$R''\subseteq R$) consist of the coordinates where $\varepsilon'$ has value strictly smaller than $k-2$ (resp.~value equal to $k-1$). Since $\varepsilon_i+1=k-1$  is odd for $i\in R$, we see that
\begin{align}
i\in R'& \implies 0\leq\left\lfloor\frac{\varepsilon_i+1}{2}\right\rfloor-\left\lfloor\frac{\varepsilon'_i+1}{2}\right\rfloor=\left\lfloor\frac{\varepsilon_i-\varepsilon'_i}{2}\right\rfloor; \label{eq: for R'}
\\
i\in R''& \implies  1=\left\lfloor\frac{\varepsilon'_i+1}{2}\right\rfloor-\left\lfloor\frac{\varepsilon_i+1}{2}\right\rfloor, \text{ as well as } \varepsilon'_i-\varepsilon_i=1.\label{eq: for R''}
\end{align}
By definition of $\varepsilon'$ we have $\sum_{i=2}^g\varepsilon_i=\sum_{i=2}^g\varepsilon_i'$, which by construction yields
\begin{equation*}
    \sum_{i\in R'}\varepsilon_i+ \sum_{i\in R''}\varepsilon_i=\sum_{i\in R'}\varepsilon_i'+\sum_{i\in R''}\varepsilon_i',
\end{equation*}
so that
\[
\sum_{i\in R'}\left(\varepsilon_i-\varepsilon_i'\right)= \sum_{i\in R''}\left(\varepsilon_i'-\varepsilon_i\right)=|R''|,
\]
where the last equality uses~\eqref{eq: for R''}. From this equality,~\eqref{eq: for R'}, and~\eqref{eq: for R''}, we get that
\begin{align*}
\sum_{i=2}^g \left(\left\lfloor\frac{\varepsilon_i'+1}{2}\right\rfloor-\left\lfloor\frac{\varepsilon_i+1}{2}\right\rfloor\right) & = \sum_{i\in R'} \left(\left\lfloor\frac{\varepsilon_i'+1}{2}\right\rfloor-\left\lfloor\frac{\varepsilon_i+1}{2}\right\rfloor\right) + \sum_{i\in R''} \left(\left\lfloor\frac{\varepsilon_i'+1}{2}\right\rfloor-\left\lfloor\frac{\varepsilon_i+1}{2}\right\rfloor\right) 
\\ 
& =\sum_{i\in R'}\left(-\left\lfloor\frac{\varepsilon_i-\varepsilon_i'}{2}\right\rfloor\right)+|R''| 
\\ 
&= \sum_{i\in R'}\left(-\left\lfloor\frac{\varepsilon_i-\varepsilon_i'}{2}\right\rfloor\right)+\sum_{i\in R'}\left(\varepsilon_i-\varepsilon_i'\right) 
\\ 
&= \sum_{i\in R'}\left(\left(\varepsilon_i-\varepsilon_i'\right)-\left\lfloor\frac{\varepsilon_i-\varepsilon_i'}{2}\right\rfloor\right)\ge 0.
\end{align*}
This verifies~\eqref{serealiza}, thereby completing the proof. 
\end{proof}

\end{document}
